\documentclass[a4paper,landscape,12pt]{article}

\usepackage{fontspec}
\usepackage{polyglossia}
\usepackage{geometry}
\usepackage{booktabs}
\usepackage{tabularx}
\usepackage{array}
\usepackage{parskip}
\usepackage{microtype}
\usepackage[table]{xcolor}
\usepackage{colortbl}

\setmainlanguage{polish}

\geometry{
  a4paper, landscape,
  top=1cm, bottom=1cm, left=1cm, right=1cm
}

\newlength{\colsep}
\setlength{\colsep}{0.8cm}
% 4:3 split of (\textwidth - \colsep)
\newlength{\leftcol}
\newlength{\rightcol}
\AtBeginDocument{
  \setlength{\leftcol}{0.44444\textwidth}
  \addtolength{\leftcol}{-0.5\colsep}
  \setlength{\rightcol}{0.55556\textwidth}
  \addtolength{\rightcol}{-0.5\colsep}
}

\definecolor{dragonred}{RGB}{160,20,20}

\setmainfont{Latin Modern Roman}
% Fallback font for card suit symbols (♥ ♦ ♠ ♣) and Ϟ
\newfontfamily\symbolfont{Segoe UI Symbol}

\newcommand{\suit}[1]{{\symbolfont #1}}
\newcommand{\drac}{{\symbolfont Ϟ}}

\setlength{\parskip}{3pt}
\setlength{\parindent}{0pt}

\newcommand{\sectiontitle}[1]{\medskip\textbf{#1}\par\smallskip}
\newcommand{\tfn}[1]{\par{\footnotesize #1}\par}
% White-text header row
\newcommand{\hdr}[1]{\color{white}\textbf{#1}}

\renewcommand{\arraystretch}{1.15}
\newcolumntype{L}[1]{>{\raggedright\arraybackslash}p{#1}}
\newcolumntype{C}[1]{>{\centering\arraybackslash}p{#1}}
\newcolumntype{R}{>{\raggedright\arraybackslash}X}

\begin{document}
\pagestyle{empty}

{\large\textbf{Gra planszowa \textit{Daredragon Fellowship} --- Ściągawka}}

\vspace{4pt}\hrule\vspace{6pt}

\noindent\begin{minipage}[t]{\leftcol}

%% ---- TABELA 1 ----
\sectiontitle{1. Obrażenia}

\rowcolors{2}{gray!15}{}
\begin{tabularx}{\linewidth}{C{0.8cm} R R R}
  \rowcolor{dragonred}
  \hdr{Obr.} & \hdr{Ataki smoka} & \hdr{Ataki obszar. smoka} & \hdr{Ataki śmiałków} \\
  13 & {[K]} Zmiażdżenie & & \\
  10 & {[D]} Ugryzienie & & \\
  8  & {[10] [9]} Stratowanie* \newline {[8]} Ciężki atak & {[A\suit{♥♦}]} Ognisty oddech & {[9]} Szarża \\
  6  & {[7] [6]} Atak & {[W]} Machnięcie ogonem* & {[8]} Precyzyjny atak \\
  5  & {[5]} Szybki atak & & {[7] [6]} Atak \newline {[5]} Szybki atak \\
  \bottomrule
\end{tabularx}
\tfn{(*) Powoduje powalenie.}

%% ---- TABELA 2 ----
\sectiontitle{2. Modyfikatory obrażeń}

\rowcolors{2}{gray!15}{}
\begin{tabularx}{\linewidth}{C{0.9cm} R R}
  \rowcolor{dragonred}
  \hdr{Poz.} & \hdr{Ataki smoka} & \hdr{Ataki śmiałków} \\
  +2\drac & Śmiałek leży na ziemi* &
    Riposta (1 tura po ataku smoka, 2 po uniku ratunkowym*) \\
  +1\drac & Atak w czasie przygotowywania akcji śmiałka* \newline \newline Więcej niż 4 graczy &
    Atak na oplątaną część ciała smoka \newline \newline
    Za każdy kolejny atak w tej samej turze tym samym kolorem \\
  -1\drac & PŻ od 1 do 13 & PŻ od 1 do 13 \\
  -2\drac & PŻ równe 0 & Śmiałek leży na ziemi \\
  -3\drac & Za każdą osłonę** & \\
  \bottomrule
\end{tabularx}
\tfn{(*) Nie dotyczy ataków obszarowych. \newline (**) Śmiałek w obszarze ataku nie osłania sam siebie.}

\end{minipage}\hspace{\colsep}\begin{minipage}[t]{\rightcol}

%% ---- TABELA 3 ----
\sectiontitle{3. Czas wykonywania akcji}

\rowcolors{2}{gray!15}{}
\begin{tabularx}{\linewidth}{C{0.85cm} R R R R}
  \rowcolor{dragonred}
  \hdr{Czas prz.} &
  \hdr{Akcje smoka, \newline czas odp. = 2} &
  \hdr{Akcje smoka, \newline czas odp. = 1} &
  \hdr{Akcje śmiałków} &
  \hdr{Akcje śmiałków} \\
  4 & {[K]} Zmiażdżenie & & {[9]} Szarża & {[Joker]} Pierwsza pomoc \\
  3 & {[D]} Ugryzienie \newline {[W]} Machnięcie ogonem & {[8]} Ciężki atak & {[8]} Precyzyjny atak & {[10]} Oplątanie \\
  2 & {[10] [9]} Stratowanie & {[7] [6]} Atak \newline {[A\suit{♥♦}]} Ognisty oddech & {[7] [6]} Atak & {[4] [3] [2]} Pomocna dłoń \\
  1 & & {[5]} Szybki atak \newline {[A\suit{♠♣}]} Przeraźliwy ryk \newline {[Joker]} Regeneracja & {[5]} Szybki atak & {[K]} Odwrócenie uwagi \\
  \bottomrule
\end{tabularx}

%% ---- TABELA 4 ----
\sectiontitle{4. Modyfikatory czasu wykonywania akcji}

\rowcolors{2}{gray!15}{}
\begin{tabularx}{\linewidth}{C{0.85cm} R R}
  \rowcolor{dragonred}
  \hdr{Czas prz.} & \hdr{Akcje smoka} & \hdr{Akcje śmiałków} \\
  +1 & PŻ równe 0 \newline \newline
       Oplątana część ciała \newline \newline
       Odwrócenie uwagi smoka w czasie przygotowywania akcji &
       Śmiałek leży na ziemi* \\
  \bottomrule
\end{tabularx}
\tfn{(*) Akcje o czasie przygotowania 4 oraz akcje typu czuwanie stają się niedostępne.}

\medskip
Czas odpoczynku smoka +1, gdy wytrącony z równowagi przez unik ratunkowy.

\end{minipage}

\end{document}
